\uuid{UbhS}
\chapitre{Probabilité discrète}
\niveau{L2}
\module{Probabilité et statistique}
\sousChapitre{Variable aléatoire discrète}
\titre{Indépendance de variables aléatoires}
\theme{variables aléatoires discrètes}
\auteur{}
\datecreate{2023-02-07}
\organisation{AMSCC}
\difficulte{2}
\contenu{

\texte{ 	Dans une boîte, il y a 6 boules rouges et 8 boules jaunes. On tire au hasard 3 boules de la boîte, sans remise. 
	Parmi les 3 boules tirées, on note $X$ le nombre de boules rouges et $Y$ le nombre de boules jaunes.}
	\begin{enumerate}
		\item \question{ Expliciter la loi de $X$. }
		\indication{Les tirages non ordonnés de $3$ boules parmi les $14$ boules sont équiprobables. Pour déterminer $\prob(X=k)$, compter les tirages contenant $k$ boules rouges et $3-k$ boules jaunes.}
		\reponse{La variable $X$ prend ses valeurs dans $\{0,1,2,3\}$. Pour tout $k\in\{0,1,2,3\}$,
		\[
		\prob(X=k)=\frac{\binom{6}{k}\binom{8}{3-k}}{\binom{14}{3}}.
		\]
		On obtient donc
		\[
		\begin{array}{|c|c|c|c|c|}
		\hline
		k & 0 & 1 & 2 & 3 \\
		\hline
		\prob(X=k) & \frac{2}{13} & \frac{6}{13} & \frac{30}{91} & \frac{5}{91} \\
		\hline
		\end{array}
		\]}
		\item \question{ Calculer l'espérance de $X$, de $Y$ et la variance de $Y$. }
		\reponse{On a
		\[
		\E(X)=\frac{9}{7}
		\quad\text{et}\quad
		\E(X^2)=\frac{207}{91}.
		\]
		Or $Y=3-X$. Par linéarité de l'espérance,
		\[
		\E(Y)=3-\E(X)=\frac{12}{7}.
		\]
		De plus, $\V(Y)=\V(3-X)=\V(X)$, donc
		\[
		\V(Y)=\E(X^2)-\E(X)^2
		=\frac{207}{91}-\left(\frac{9}{7}\right)^2
		=\frac{396}{637}.
		\]}
		\item \question{ Les variables $X$ et $Y$ sont-elles indépendantes ? Justifier. }
		\reponse{Les variables ne sont pas indépendantes : elles vérifient toujours $X+Y=3$. Par exemple,
		\[
		\prob(X=0,Y=3)=\prob(X=0)=\frac{2}{13},
		\]
		alors que
		\[
		\prob(X=0)\prob(Y=3)=\frac{2}{13}\times\frac{2}{13}=\frac{4}{169}.
		\]
		Ces deux probabilités sont différentes.}
	\end{enumerate}}
